% Part: first-order-logic
% Chapter: natural-deduction
% Section: soundness-identity

\documentclass[../../../include/open-logic-section]{subfiles}

\begin{document}

\olfileid[fr]{fol}{ntd}{sid}

\olsection{Correction avec égalité}

\begin{prop}
La déduction naturelle munie des règles pour $\eq$ est correcte.
\end{prop}

\begin{proof}
Toute !!{formula} de la forme $\eq[t][t]$ est valide : pour toute
!!{structure}~$\Struct M$, on a $\Sat{M}{\eq[t][t]}$.
Le terme $t$ est supposé clos, c'est-à-dire sans variable;
les affectations de variables n'interviennent donc pas.

Supposons que la dernière inférence d'une !!{derivation} soit
\Elim{\eq}. Considérons d'abord l'orientation suivante :
\begin{prooftree}
  \AxiomC{$\Gamma_1$}
  \RightLabel{$\delta_1$}
  \DeduceC{$\eq[t_1][t_2]$}
  \AxiomC{$\Gamma_2$}
  \RightLabel{$\delta_2$}
  \DeduceC{$!A(t_1)$}
  \RightLabel{\Elim{\eq}}
  \BinaryInfC{$!A(t_2)$}
\end{prooftree}
Les prémisses $\eq[t_1][t_2]$ et $!A(t_1)$ sont dérivées à
partir des hypothèses !!{undischarged}s~$\Gamma_1$ et $\Gamma_2$,
respectivement. Il faut montrer que $!A(t_2)$ est une conséquence
de $\Gamma_1 \cup \Gamma_2$.
Considérons une !!{structure}~$\Struct{M}$ telle que
$\Sat{M}{\Gamma_1 \cup \Gamma_2}$.
Par hypothèse de récurrence, $\Sat{M}{!A(t_1)}$ et
$\Sat{M}{\eq[t_1][t_2]}$. Donc
$\Value{t_1}{M} = \Value{t_2}{M}$.
Soit $s$ une affectation de variables quelconque, et posons
$m = \Value{t_1}{M} = \Value{t_2}{M}$.
Par la \olref[fol][syn][ext]{prop:ext-formulas},
$\Sat{M}{!A(t_1)}[s]$ si et seulement si
$\Sat{M}{!A(x)}[\Subst{s}{m}{x}]$, si et seulement si
$\Sat{M}{!A(t_2)}[s]$.
Puisque $\Sat{M}{!A(t_1)}$, il vient $\Sat{M}{!A(t_2)}$.
L'autre orientation de \Elim{\eq} se traite de la même façon,
en échangeant les rôles de $t_1$ et $t_2$; l'égalité de leurs
dénotations est symétrique.
\end{proof}

\begin{editorial}
La source détaille une seule orientation de l'élimination de
l'égalité. Le traitement analogue de la seconde est explicité
pour couvrir les deux règles présentées dans la section précédente.
\end{editorial}

\end{document}
